% EGA V, English translation by P. Blass, 1987 (cote 24): archive cote 24, site page 52 (PDF page 53).
% typescript. Draft compared: none.
% AI-transcribed by Claude Opus 5.5 (claude-opus-5-5, adaptive thinking, effort high), Message Batches API, batch ega24b (msgbatch_01RGJJeHdR96U16ZWFuUjtDx), 2026-11-13; not yet checked by eye.
\scanpage{24/52}
\mspage{52}

non-osculating hyperplane for X at a point x $\in$ X(k) such that
there exists another point x$^{1}$ $\in$ X(k) at which the same hyperplane
is tangent at x [or x$^{1}$ illegible]*. Indeed there are obvious
examples with $\underline{P}$ of dim two, X a non-singular curve of degree $\geq$ 4,
in any characteristic. [Note here: the ``double tangents'' of X
correspond to the double points of the ``dual curve.'']

\uline{Corollary 8.11}

Let us assume that k has characteristic zero. Then (a) The
image of $\underline{Y}$sup sing in $\underline{P}$v is of codimension $\geq$ 2. (b) the condition
(iii) of 8.8 is equivalent to other conditions i.e. the negation
of the other conditions, let us assume that $\underline{Y}$sing = Ysup sing], means
also that the image of $\underline{Y}$sing [or Ysup sing illegible]* in $\underline{P}$v is of
codimension $\geq$ 2.

Evidently, the assertion (b) implies (a) taking into account
8.8. But by dimension theory, (iii) means that L/K is a finite
extension (we could put it in the form ((iii) bis) in 8.8 and in
characteristic zero L is always separable over K hence the
condition (i bis) of 8.8.

\uline{Remark 8.12}

Geometrically the assertion (a) means essentially that for a
sufficiently general linear pencil of hyperplane sections every
member of the pencil is smooth or has for s ??? [illegible]* only
for geometric points singular points ordinary double points (and
in fact as one sees immediately it can be said in statement (a)
consequently in a form a little more precise--we have at most \uline{one}

\margin{* ASK AG (or figure out)}
